\section{One Hida ideal, the marked theorem and Case V}

% BEGIN M72S1
\paragraph{M72S1. Filling the precise arithmetic hypothesis of the Hida construction.}
Form the rank-general literature hypothesis with V72L8 for its local
entry, A72F9 for the automorphic entry, F71A3 for Galois cohomology,
and C72M1--11 for the eleven CM and abelian-scheme entries. The
divided-power comparisons used there are D72P1--3 and apply over
every base ring in the original definitions. In particular the
prime bound in the tame Galois comparison is used only in its
specified branch, and the Hodge filtration is the one on the
given abelian scheme.

For the remaining fields of $H'=\mathrm{Af.HypothesesH}'$, take
O72H1 for $\mathrm{cohC}$ and $\mathrm{conjC}$, O72H2 for
$\mathrm{borelC}$, O72H6 for $\mathrm{galR}$, and O72H3 for
$\mathrm{smoR},\mathrm{gkR},\mathrm{admR}$. Take O72H4 for the
quadratic entry $\mathrm{bc}$ and O72H5 for $\mathrm{newform}$.
This fills the parent record. The five additional fields are the
weak descent of F71A2, the two cyclic assertions of O72H4, the
definite quaternion algebra of F71A1, and cuspidal semisimplicity
from O72H3. These are the entries of the original structure.
The descent of a cusp to the definite unitary group and the
base change of a possibly nondiscrete-series local factor out of
that group have different conclusions; the latter occurs inside
A72F6, in the rank-general record. Both directions have now been
supplied.
% END M72S1

% BEGIN M72S2
\paragraph{M72S2. An ideal and the automorphy system on that same ideal.}
Apply HIA2 with M72S1. On each admissible setup form the module
of ordinary automorphic forms and its actual Hecke algebra at the
empty auxiliary level, over the given coefficient DVR. In the
nonzero case the map from the universal ring $B$ onto the Hecke
algebra gives $J$ as its kernel, equivalently the annihilator of
the module. The zero-module case uses the unit ideal and the
point implication of HIA1. HIA1 supplies all six Hida properties,
and the auxiliary levels of HIA2 supply $\mathrm{HasAut}(B,J)$
for that very $J$. This proves $\mathrm{HidaUAlt}$ under the
original setup hypotheses. In particular it has the ideal before
any arithmetic point or chain is tested.

A71M1 then applies the two prime tests on this ideal. G71H4 and
KX3 supply the cohomology and local reciprocity used in HIA3.
Consequently $\mathrm{ModUAlt}$ holds, and L69D3 supplies the
finite assertion and both rank-three and rank-two automorphy
assertions on every eligible original link.
% END M72S2

% BEGIN M72S3
\paragraph{M72S3. The whole original marked theorem.}
Fix the original Gaussian host data, the embedding $\iota$ and an
eligible frame. Apply A71M2 with the ideal supplier of M72S2.
Concretely, AFE16 gives the marked arithmetic frame; the finite
test of M72S2 returns on it by AFE20. The rank-three test there
uses the twist $\varepsilon_L^{-1}$, and D70W5 carries its
automorphy and the specified finite-place bounds back to the
given Gaussian field. M65A4 and the marked comparison give
one $g_1$ in the host quotient with the stated marking identities.
For every universal pair $(R,u)$ at that mark, the weight-zero
part is finite over the coefficient DVR. For every allowed entry
map and every type-I point of that pair, the associated matrices
give $\mathrm{AutOnAlt}(K,\iota,-1)$ with the prescribed twist.
The same $g_1$ works for both assertions and for every universal
pair. These are exactly the two conclusions of
$\mathrm{ThmS6}_s$ with its original premises. This discharges
the whole obligation \texttt{THM\_S6\_S} in the human reconstruction.
% END M72S3

% BEGIN M72S4
\paragraph{M72S4. The Case V estimate with its original guard.}
Use the marked conclusion M72S3 in HVJ1 on the Case V frame under
the original inequality $R_1\le2^{c_2(\psi)}$. The previous host
vanishing, rigidity and frame-length comparisons, the two algebraic
lemmas and the cohomological inputs G71H1--3 give the finite strict
Selmer group and the scaled length bounded by a sum of trace
differences on that frame. HVJ2 uses the uniform bound VUN6 for
those same points and their common Frobenius element. With
$X=\ell N e\ge2$, its final estimate is
$\mathrm{Vn}\log\ell\le 2X^C\le X^{C+1}$, with a constant
fixed before the target character and the host. This supplies the
dependent Case V producer in the register. It counts once as a
consequence of closing the marked theorem.
% END M72S4
